%% rsuserguide.tex --- user guide and reference manual

% Copyright (C) 2023 Ralph Schleicher

% Redistribution and use in source and binary forms, with or without
% modification, are permitted provided that the following conditions
% are met:
%
%    * Redistributions of source code must retain the above copyright
%      notice, this list of conditions and the following disclaimer.
%
%    * Redistributions in binary form must reproduce the above copyright
%      notice, this list of conditions and the following disclaimer in
%      the documentation and/or other materials provided with the
%      distribution.
%
%    * Neither the name of the copyright holder nor the names of its
%      contributors may be used to endorse or promote products derived
%      from this software without specific prior written permission.
%
% THIS SOFTWARE IS PROVIDED BY THE COPYRIGHT HOLDERS AND CONTRIBUTORS
% "AS IS" AND ANY EXPRESS OR IMPLIED WARRANTIES, INCLUDING, BUT NOT
% LIMITED TO, THE IMPLIED WARRANTIES OF MERCHANTABILITY AND FITNESS
% FOR A PARTICULAR PURPOSE ARE DISCLAIMED.  IN NO EVENT SHALL THE
% COPYRIGHT HOLDER OR CONTRIBUTORS BE LIABLE FOR ANY DIRECT, INDIRECT,
% INCIDENTAL, SPECIAL, EXEMPLARY, OR CONSEQUENTIAL DAMAGES (INCLUDING,
% BUT NOT LIMITED TO, PROCUREMENT OF SUBSTITUTE GOODS OR SERVICES;
% LOSS OF USE, DATA, OR PROFITS; OR BUSINESS INTERRUPTION) HOWEVER
% CAUSED AND ON ANY THEORY OF LIABILITY, WHETHER IN CONTRACT, STRICT
% LIABILITY, OR TORT (INCLUDING NEGLIGENCE OR OTHERWISE) ARISING IN
% ANY WAY OUT OF THE USE OF THIS SOFTWARE, EVEN IF ADVISED OF THE
% POSSIBILITY OF SUCH DAMAGE.

%% Code:

\documentclass{scrartcl}
\usepackage{rspackages}
\usepackage{amssymb} %for \checkmark
\usepackage{scrlogo} %for \KOMAScript
\usepackage{rsmacros}

\makeatletter

\NewDocumentCommand\env{m}{\code{#1}\thinspace\textsubscript{env}}
\NewDocumentCommand\sty{m}{\code{#1}\thinspace\textsubscript{sty}}
\NewDocumentCommand\cls{m}{\code{#1}\thinspace\textsubscript{cls}}

% An index for the symbol names, i.e. LaTeX control sequences and
% environments.
\newindex[Symbol Index]{sym}

% Add an entry to the symbol index.
%
% The mandatory argument is the symbol name.
% The optional argument provides the sorting rule.
% The star variant makes no assumption about the formatting of the
%  symbol name.  Default is to format it with the ‘\symb’ command.
\NewDocumentCommand\symindex{s o m}{%
  \IfBooleanTF{#1}{%
    % Star is present.
    \IfNoValueTF{#2}{%
      \index[sym]{#3}}{%
      \index[sym]{#2@#3}}}{%
    % Without star.
    \IfNoValueTF{#2}{%
      \index[sym]{#3@\symb{#3}}}{%
      \index[sym]{#2@\symb{#3}}}}}

% TeX control sequence.
\NewDocumentCommand\texindex{m}{%
  \symindex*{#1@\code{\\#1}}}
\NewDocumentCommand\texindexvisiblespace{}{%
  \symindex*{\ @\code{\\\ }}}
% LaTeX environment.
\NewDocumentCommand\envindex{m}{%
  \symindex*{#1@\env{#1}}}
% LaTeX package.
\NewDocumentCommand\styindex{m}{%
  \index{#1@\sty{#1}}}

\NewDocumentCommand\anchor{m}{%
  \hypertarget{#1}{}}
\NewDocumentCommand\texanchor{m}{%
  \hypertarget{tex:#1}{}}
\NewDocumentCommand\envanchor{m}{%
  \hypertarget{env:#1}{}}
\NewDocumentCommand\texref{m}{%
  \hyperlink{tex:#1}{\code{\\#1}}}
\NewDocumentCommand\envref{m}{%
  \hyperlink{env:#1}{\code{#1}}}

\makeatother

\title{Ralph's \LaTeX{} Macros}
\author{Copyright \copyright{} 2023 Ralph Schleicher\footnote{%
Permission is granted to copy, distribute, and/or modify this
document under the terms of the GNU Free Documentation License,
Version 1.3 or any later version published by the Free Software
Foundation; with no Invariant Sections, no Front-Cover Texts,
and no Back-Cover Texts.}}
\date{Version 1.1}

\begin{document}
\maketitle
\tableofcontents


\section{Introduction}

Ralph's \LaTeX{} macros is a collection of frequently used packages,
commands, and environments for writing \LaTeX{} documents.  The first
part is a selection of useful standard \LaTeX{} packages.  The second
part are Ralph's own \LaTeX{} commands and environments, which exist
either because there is no standard \LaTeX{} equivalent or as a matter
of style.


\section{Standard \LaTeX{} Packages}

\styindex{rspackages}
The \code{rspackages} package loads a selection of standard \LaTeX{}
packages.  That is,
\begin{description}
\item[\href{https://mirrors.ctan.org/macros/latex/base/encguide.pdf}{\code{fontenc}}]
\styindex{fontenc}
\index{font encoding}
\index{encoding!font}
If the \TeX{} engine is \code{latex} or \code{pdflatex}, enable the
T1 font encoding.  Otherwise, do nothing.

\item[\href{https://mirrors.ctan.org/macros/latex/base/ifthen.pdf}{\code{ifthen}}]
\styindex{ifthen}
\texindex{ifthenelse}
\index{conditional}
This package provides the \symb{\ifthenelse} command.

\item[\href{https://mirrors.ctan.org/macros/latex/base/alltt.pdf}{\code{alltt}}]
\styindex{alltt}
\envindex{alltt}
\envindex{verbatim}
This package provides the \code{alltt} environment, which is like the
\code{verbatim} environment, except that the escape character \code{\\}
and the grouping characters \code{\{} and \code{\}} have their usual
meaning.

\item[\href{https://mirrors.ctan.org/macros/latex/required/tools/array.pdf}{\code{array}}]
\styindex{array}
\envindex{array}
\envindex{tabular}
This package provides the \code{array} and \code{tabular} environments.

\item[\href{https://mirrors.ctan.org/macros/latex/required/tools/dcolumn.pdf}{\code{dcolumn}}]
\styindex{dcolumn}
\index{decimal point}
\index{alignment!decimal point}
This package provides a system for aligning table column entries at
the decimal point.

\item[\href{https://mirrors.ctan.org/macros/latex/required/tools/longtable.pdf}{\code{longtable}}]
\styindex{longtable}
\envindex{longtable}
\index{multi-page table}
\index{table!multi-page}
This package provides the \code{longtable} environment for typesetting
multi-page tables.

\item[\href{https://mirrors.ctan.org/macros/latex/required/tools/tabularx.pdf}{\code{tabularx}}]
\styindex{tabularx}
\envindex{tabularx}
\envindex{tabular*}
This package provides the \code{tabularx} environment, which is like
the \code{tabular*} environment, but modifies the widths of certain
columns, rather than the inter column space, to set a table with the
requested total width.

\item[\href{https://mirrors.ctan.org/macros/latex/required/tools/multicol.pdf}{\code{multicol}}]
\styindex{multicol}
\envindex{multicols}
\index{multi-column format}
\index{format!multi-column}
This package provides the \code{multicols} environment which allows
switching between single-column and multi-column format on the same
page.

\item[\href{https://mirrors.ctan.org/macros/latex/required/tools/verbatim.pdf}{\code{verbatim}}]
\styindex{verbatim}
\envindex{verbatim}
\texindex{verbatiminput}
This package provides a new implementation of the \code{verbatim}
environment.  It also defines the \symb{\verbatiminput} command to
input a whole file verbatim.

\item[\href{https://mirrors.ctan.org/macros/latex/contrib/ragged2e/ragged2e.pdf}{\code{ragged2e}}]
\styindex{ragged2e}
\envindex{center}
\envindex{flushleft}
\envindex{flushright}
\texindex{centering}
\texindex{raggedleft}
\texindex{raggedright}
This package provides a new implementation of the text alignment
environments \code{center}, \code{flushleft}, and \code{flushright},
and the respective \symb{\centering}, \symb{\raggedleft}, and
\symb{\raggedright} commands.

\item[\href{https://mirrors.ctan.org/macros/latex/contrib/enumitem/enumitem.pdf}{\code{enumitem}}]
\styindex{enumitem}
\envindex{enumerate}
\envindex{itemize}
\envindex{description}
This package provides a new implementation of the three basic list
environments \code{enumerate}, \code{itemize}, and \code{description}.

\item[\href{https://mirrors.ctan.org/macros/latex/required/amsmath/amsldoc.pdf}{\code{amsmath}}]
\styindex{amsmath}
\index{alignment!equation}
This package provides miscellaneous enhancements for typesetting
mathematical formulas.  Displayed equations are aligned to the left
margin and equation numbers are placed on the right.  For integral
symbols, large operator symbols, and named operators, limits are
typeset below and above the symbol.

\item[\href{https://mirrors.ctan.org/macros/latex/contrib/siunitx/siunitx.pdf}{\code{siunitx}}]
\styindex{siunitx}
\index{physical quantity}
\index{quantity}
\index{number}
\index{unit}
This package provides commands for typesetting physical quantities and
units in a consistent way.

\item[\href{https://mirrors.ctan.org/macros/latex/contrib/xcolor/xcolor.pdf}{\code{xcolor}}]
\styindex{xcolor}
\styindex{color}
\index{color}
This package provides extensions to the basic facilities of the
\code{color} package.

\item[\href{https://mirrors.ctan.org/macros/latex/required/graphics/grfguide.pdf}{\code{graphicx}}]
\styindex{graphicx}
\styindex{graphics}
\texindex{includegraphics}
This package provides extensions to the basic facilities of the
\code{graphics} package.  It also defines the \symb{\includegraphics}
command.

\item[\href{https://mirrors.ctan.org/macros/latex/contrib/pdflscape/pdflscape.pdf}{\code{pdflscape}}]
\styindex{pdflscape}
\envindex{landscape}
This package provides PDF support for the \code{landscape} environment.

\item[\href{https://mirrors.ctan.org/macros/latex/contrib/splitindex/splitidx.pdf}{\code{splitidx}}]
\styindex{splitidx}
\texindex{newindex}
\texindex{index}
This package provides support for creating multiple indices.  The
\symb{\newindex} command can be used to create an additional index.
The \symb{\index} command is redefined so that you can add an index
entry to a dedicated index.  The \code{splitidx} package is loaded
with the \code{protect} option to avoid premature expansion of
\symb{\index} arguments.

\item[\href{https://mirrors.ctan.org/macros/latex/contrib/hyperref/doc/hyperref-doc.pdf}{\code{hyperref}}]
\styindex{hyperref}
\index{hypertext link}
\index{link!hypertext}
This package provides extensive support for hypertext links in
\LaTeX{}.  Links are neither framed nor coloured.
\end{description}
The \code{rspackages} package is best loaded after the page layout and
document font is defined, e.g.
\begin{display}
\begin{verbatim}
\documentclass{scrartcl}
\usepackage{cmbright}
\usepackage{rspackages}
\end{verbatim}
\end{display}


\subsection{Layout and Formatting}

All the loaded packages are configured for a consistent style.

\begin{definition}
\texindex{rs@indent}
\texanchor{rs@indent}
\defdimen{\rs@indent}
The basic indentation for lists and displays.
Default is \qty{10}{mm}.
\end{definition}

\index{nested list}
\index{list!nested}
Nested lists have the same relative indentation.  The default label
for itemized lists is an en-dash.  Description lists don't make any
assumption about the formatting of labels.

Sometimes you may want to typeset some content in a slightly smaller
font so that the whole thing does not take up too much space.  There
is a command for this purpose which sets the font size accordingly.

\begin{definition}
\texindex{smallsize}
\texanchor{smallsize}
\deftex{\smallsize}
The font size for small displays.
Default is \symb{\footnotesize}.
\end{definition}


\subsection{Managing Indices}

\index{index}
\texindex{makeindex}
\texindex{printindex}
\envindex{theindex}
The \code{rspackages} package calls \symb{\makeindex} and it provides
a \symb{\printindex} command that does the right thing, whether or not
you use a single index or multiple indices.  The \code{theindex}
environment is redefined so that it only prints the actual index in
double-column format.  That means, you must provide the sectioning
command when printing an index in your document.  Here is an example
for the indices of a programming reference manual.
\begin{display}
\begin{verbatim}
\chapter{Index}
\section*{Concept Index}
\printindex
\section*{Symbol Index}
\printindex[sym]
\end{verbatim}
\end{display}
This behaviour of the \code{theindex} environment gives you full
control over how to print an index, whether or not it appears in the
table of contents, and what is printed before and after it -- without
further ado.

\index{index!style file}
\index{style file!index}
\index{rspackages.ist@\code{rspackages.ist}}
\index{jobname.ind@\code{\meta{job-name}.ind}}
An index style file ‘\code{rspackages.ist}’ is installed along with
the \code{rspackages} package.  The recommended shell command to
create the index file \code{\meta{job-name}.ind} is as follows.
\index{makeindex@\code{makeindex}}
\begin{display}
\begin{alltt}
makeindex -c -s rspackages.ist \meta{job-name}.idx
\end{alltt}
\end{display}
\index{splitindex@\code{splitindex}}
If you use multiple indices, call \code{splitindex} instead.
\begin{display}
\begin{alltt}
splitindex -v \meta{job-name}.idx -- -c -s rspackages.ist
\end{alltt}
\end{display}
When the index style file ‘\code{rspackages.ist}’ is applied, the
final index files have the following format.
\begin{display}
\begin{verbatim}
\begin{theindex}
\indexgroup{\indexsymbolsname}
\item \$ (dollar sign)\indexpagesep 1
\indexspace
\indexgroup{\indexnumbersname}
\item 0 (zero)\indexpagesep 1
\indexspace
\indexgroup{A}
\item an entry\indexpagesep 1
\end{theindex}
\end{verbatim}
\end{display}

\begin{definition}
\texindex{indexsymbolsname}
\texanchor{indexsymbolsname}
\deftex{\indexsymbolsname}
The name for the symbols group.
Default is “\indexsymbolsname”.
\end{definition}

\begin{definition}
\texindex{indexnumbersname}
\texanchor{indexnumbersname}
\deftex{\indexnumbersname}
The name for the numbers group.
Default is “\indexnumbersname”.
\end{definition}

\begin{definition}
\texindex{indexspace}
\texanchor{indexspace}
\deftex{\indexspace}
The space between adjacent groups.
Default is \symb{\medskip}.
\end{definition}

\begin{definition}
\texindex{indexgroup}
\texanchor{indexgroup}
\deftex{\indexgroup} \meta{name}
The sectioning command for the group header.
Default is \symb{\minisec}.

Argument \meta{name} is the header of the group.
\end{definition}

\begin{definition}
\texindex{indexpagesep}
\texanchor{indexpagesep}
\deftex{\indexpagesep}
The delimiter between the index entry and the page number.
Default is an em-space.
\end{definition}

The default behaviour of the \code{theindex} environment is to omit
group headers and only insert \texref{indexspace} between groups.  If
you redefine the \code{theindex} environment, you must call
\texref{indexsetup} at the beginning of the index.
\begin{display}
\begin{verbatim}
\RenewDocumentEnvironment{theindex}{}{%
  % The ‘multicols*’ environment does not balance
  % the columns on the last page.
  \begin{multicols*}{2}
  \indexsetup[headers]}{%
  \end{multicols*}}
\end{verbatim}
\end{display}

\begin{definition}
\texindex{indexsetup}
\texanchor{indexsetup}
\deftex{\indexsetup} \meta{options}\kleene?
Adjust the typesetting commands and formatting parameters for the
index.  Must be called locally in a \TeX{} group.

Optional argument \meta{options} is a comma separated list of
options to change the formatting of the index.
\begin{description}
\item[\code{headers}]
Print group headers.

\item[\code{nogroups}]
Disable grouping, i.e. omit group headers and disable
\symb{\indexspace}.

\item[\code{comma}]
Separate the index entry and the page number by a comma followed
by a space.
\end{description}
\end{definition}


\section{Ralph's \LaTeX{} Macros}

\styindex{rsmacros}
Ralph's \LaTeX{} macros can be loaded all at once by using the
\code{rsmacros} package or individually as documented in the
respective section.


\subsection{Displayed Text}

\styindex{rsdisplay}
The \code{rsdisplay} package provides the \code{display} environment
for displayed text.

% display
% ----------------------------------------------------------------------
\begin{definition}
\envindex{display}
\envanchor{display}
\defenv{display} \code{*}\kleene? \meta{options}\kleene?
Like the \code{displaymath} environment but for vertical material.
This environment indents the left margin and inserts space above and
below the text body.
\begin{itemize}
\item The optional star argument omits the indentation.  This is a
nice alternative for a ‘here’ float.
\item Optional argument \meta{options} is a comma separated list of
options.
\begin{description}
\codeitem{big}
Vertical space around the displayed content is a
\symb{\bigskipamount}.  This is the default.
\codeitem{medium}
Vertical space around the displayed content is a
\symb{\medskipamount}.
\codeitem{small}
Vertical space around the displayed content is a
\symb{\smallskipamount}.
\codeitem{par}
Vertical space around the displayed content is a
\symb{\parskip}.
\codeitem{none}
No vertical space around the displayed content.
\end{description}
\end{itemize}
\end{definition}

For example,
\begin{display}
\begin{verbatim}
\begin{display}
Some text to be displayed.
\end{display}
\end{verbatim}
\end{display}
is printed as
\begin{display}
Some text to be displayed.
\end{display}
Displayed math
\begin{displaymath}
\sin^2 x + \cos^2 x = 1
\end{displaymath}
has the same spacing.


\subsection{Code and Syntax}

\styindex{rsmetasyntax}
The \code{rsmetasyntax} package provides macros for indicating
commands, symbol names (identifiers), and meta-syntactic variables.

% code
% ----------------------------------------------------------------------
\begin{definition}
\texindex{code}
\texanchor{code}
\deftex{\code} \code{*}\kleene? \meta{text}
\index{literal text}
Print \meta{text} as literal text in a typewriter type font.
The special characters
\begin{display}
\code{\\}\quad
\code{\{}\quad
\code{\}}\quad
\code{\#}\quad
\code{\$}\quad
\code{\%}\quad
\code{\&}\quad
\code{\^}\quad
\code{\_}\quad
\code{\~}
\end{display}
must be escaped.
\index{backslash}
\symindex{\\}
\index{visible space}
\texindexvisiblespace
Two control sequences have another meaning than usual; \code{\\\\}
prints a backslash character and \code{\\\ } prints a visible space.
The star variant allows line breaks in \meta{text}.  Whitespace is not
significant in \meta{text} but the active tilde character can be used
to inserts a no-break space.

As a rule of thumb, single word commands, options, and symbol names
should be printed without quotes.  Examples on the other hand, should
be printed with single quotes.
\end{definition}

% symb
% ----------------------------------------------------------------------
\begin{definition}
\texindex{symb}
\texanchor{symb}
\deftex{\symb} \meta{name}
\index{symbol name}
\index{identifier}
Print \meta{name} as a symbol, i.e. as a command or programming
language identifier.  The formatting of the \symb{\symb} macro is like
\symb{\code} but all characters in \meta{name} print itself.  For
example, ‘\code{\\symb\{\\symb\}}’ prints \symb{\symb}.

Please note that the grouping characters \code{\{} and \code{\}} have
to be nested properly in \meta{name} and that multiple consecutive
space characters are collapsed into a single space character.
Additionally, all leading and trailing space characters are ignored.
\end{definition}

\texindex{verb}
Contrary to the standard \LaTeX{} \symb{\verb} command, the
\symb{\code} and \symb{\symb} command can be used inside an
argument of a \TeX{} control sequence.  In a list environment,
for example, ‘\code{\\item[\\symb\{name\}]}’ works as expected.

\begin{definition}
\texindex{codeitem}
\texanchor{codeitem}
\deftex{\codeitem} \meta{text}
The \symb{\codeitem} command is syntactic sugar for
‘\code{\\item[\\code\{\meta{text}\}]}’.
\end{definition}

\begin{definition}
\texindex{symbitem}
\texanchor{symbitem}
\deftex{\symbitem} \meta{name}
The \symb{\symbitem} command is syntactic sugar for
‘\code{\\item[\\symb\{\meta{name}\}]}’.
\end{definition}

\index{meta-syntactic variable}
\index{variable!meta-syntactic}
Sooner or later you want to use placeholders inside of code snippets.
The following commands print such meta-syntactic variables.

% var
% ----------------------------------------------------------------------
\begin{definition}
\texindex{var}
\texanchor{var}
\deftex{\var} \meta{name}
\index{meta-syntactic variable}
\index{variable!meta-syntactic}
Print \meta{name} as a meta-syntactic variable.  The formatting of the
\symb{\var} macro adopts the style of variables in mathematics.  For
example, ‘\code{\\var\{name\}}’ prints \var{name}.
\end{definition}

% meta
% ----------------------------------------------------------------------
\begin{definition}
\texindex{meta}
\texanchor{meta}
\deftex{\meta} \meta{name}
\index{meta-syntactic variable}
\index{variable!meta-syntactic}
\index{non-terminal symbol}
\index{symbol!non-terminal}
Print \meta{name} as a meta-syntactic variable delimited by angle
brackets.  The formatting of the \symb{\meta} macro adopts the style
of a non-terminal symbol in Backus–Naur form notation.  For example,
‘\code{\\meta\{name\}}’ prints \meta{name}.
\end{definition}

\index{occurrence}
\index{repetition}
Operators are frequently used to describe repeating syntactic
elements.  The Kleene star is quite commonly used to denote zero
or more occurrences.  Likewise the Kleene plus for one or more
occurrences.  In regular expressions the question mark denotes
an optional pattern, i.e. zero or one occurrence.

% kleene
% ----------------------------------------------------------------------
\begin{definition}
\texindex{kleene}
\texanchor{kleene}
\deftex{\kleene} \meta{symbol}
\index{Kleene operator}
Produce a Kleene operator for marking repetitive elements of a
grammar; \code{\\kleene*} and \code{\\kleene+} produce the Kleene
star and Kleene plus unary operator respectively.  For example,
‘\code{\\meta\{expr\}\\kleene*}’ prints \meta{expr}\kleene*, which
means zero or more occurrences of \meta{expr}.
\end{definition}

% metasyntax
% ----------------------------------------------------------------------
\begin{definition}
\envindex{metasyntax}
\envanchor{metasyntax}
\defenv{metasyntax}
\index{syntax}
\index{grammar}
An environment to describe the syntax of a language, for example, a
computer programming language, an application programming interface, a
file or document format, or a communication protocol.

Regular text is printed as in the \symb{\code} command and such texts
denote terminals.  Non-terminals should be formatted with the
\symb{\meta} command.  A set of meta-symbol commands is available
inside the body of the \symb{metasyntax} environment denoting syntax
contructs.
\begin{description}
\item[\code{\\=}]
\symindex{\=}
\index{syntax!production}
Production or definition.
\begin{multicols}{2}
\begin{verbatim}
\meta{foo}\=\meta{bar}
\end{verbatim}
\columnbreak
\begin{metasyntax}
\meta{foo}\=\meta{bar}
\end{metasyntax}
\end{multicols}

\item[\code{\\,}]
\symindex{\,}
\index{syntax!sequence}
Sequence or concatenation.
\begin{multicols}{2}
\begin{verbatim}
\meta{foo}\,\meta{bar}
\end{verbatim}
\columnbreak
\begin{metasyntax}
\meta{foo}\,\meta{bar}
\end{metasyntax}
\end{multicols}

\item[\code{\\|}]
\symindex{\|}
\index{syntax!alternative}
Alternative or choice.
\begin{multicols}{2}
\begin{verbatim}
\meta{foo}\|\meta{bar}
\end{verbatim}
\columnbreak
\begin{metasyntax}
\meta{foo}\|\meta{bar}
\end{metasyntax}
\end{multicols}

\item[\code{\\*}]
\symindex{\*}
\index{syntax!repetition}
\index{Kleene star operator}
\index{operator!Kleene star}
Repetition, zero or more times (Kleene star).
\begin{multicols}{2}
\begin{verbatim}
\meta{foo}\*
\end{verbatim}
\columnbreak
\begin{metasyntax}
\meta{foo}\*
\end{metasyntax}
\end{multicols}

\item[\code{\\+}]
\symindex{\+}
\index{syntax!repetition}
\index{Kleene plus operator}
\index{operator!Kleene plus}
Repetition, one or more times (Kleene plus).
\begin{multicols}{2}
\begin{verbatim}
\meta{foo}\+
\end{verbatim}
\columnbreak
\begin{metasyntax}
\meta{foo}\+
\end{metasyntax}
\end{multicols}

\item[\code{\\?}]
\symindex{\?}
\index{syntax!optional}
Optional phrase.
\begin{multicols}{2}
\begin{verbatim}
\meta{foo}\?
\end{verbatim}
\columnbreak
\begin{metasyntax}
\meta{foo}\?
\end{metasyntax}
\end{multicols}

\item[\code{\\(...\\)} or \code{\\group\{...\}}]
\texindex{group}
\index{syntax!group}
Grouping.
\begin{multicols}{2}
\begin{verbatim}
\(foo\|bar\)
\group{foo\|bar}
\end{verbatim}
\columnbreak
\begin{metasyntax}
\(foo\|bar\)
\group{foo\|bar}
\end{metasyntax}
\end{multicols}

\item[\code{\\start}]
\texindex{start}
\index{syntax!start rule}
Mark the start rule.
\begin{multicols}{2}
\begin{verbatim}
\start\meta{foo}
\end{verbatim}
\columnbreak
\begin{metasyntax}
\start\meta{foo}
\end{metasyntax}
\end{multicols}

\item[\code{\\default}]
\texindex{default}
\index{syntax!default value}
Mark the default value.
\begin{multicols}{2}
\begin{verbatim}
\default{42}
\end{verbatim}
\columnbreak
\begin{metasyntax}
\default{42}
\end{metasyntax}
\end{multicols}

\item[\code{\\\&}]
\symindex{\&}
Continuation, add it at the beginning of a line.
\begin{multicols}{2}
\begin{verbatim}
\meta{list}\=\meta{item}
\&\|\meta{item}\,\meta{list}
\end{verbatim}
\columnbreak
\begin{metasyntax}
\meta{list}\=\meta{item}
\&\|\meta{item}\,\meta{list}
\end{metasyntax}
\end{multicols}
Each \symb{\&} increases the indentation by one unit.

\item[\code{\\\\}, \code{\\\ }, \code{\\dots}]
\index{backslash}
\symindex{\\}
\index{visible space}
\texindexvisiblespace
\index{ellipses}
\texindex{dots}
Backslash, visible space, and ellipses.
\begin{multicols}{2}
\begin{metasyntax}
\\\\
\\\ {}
\\dots
\end{metasyntax}
\columnbreak
\begin{metasyntax}
\\
\ {}
\dots
\end{metasyntax}
\end{multicols}
\end{description}
\end{definition}


\subsection{Documenting Definitions}

\styindex{rsdefinition}
\index{documentation}
The \code{rsdefinition} package provides macros to describe functions,
variables, data types, and similar entities in a consistent style.
The interface is modeled after the respective Texinfo functionality.

For example,
\begin{display}
\begin{verbatim}
\begin{definition}
\defun{forward-line} &optional \var{n}
Move \var{n} lines forward (backward if \var{n} is negative).
\end{definition}
\end{verbatim}
\end{display}
is printed as
\begin{display}
\begin{definition}
\defun{forward-line} &optional \var{n}
Move \var{n} lines forward (backward if \var{n} is negative).
\end{definition}
\end{display}

Start the definition with a \symb{\def}\ldots{} command to insert
the signature of the entity.  The body of the definition, i.e. the
documentation, follows on succeeding lines.

The signature line is printed in the \texref{definitionfont} font.
The name of the entity is formatted using the \texref{symb} command.
If the entity has parameters, the entire parameter list is either a
balanced group or the rest of the line.

The parameter list has several features to format Lisp lambda lists
nicely.  Parentheses, \code( and \code), are active and print itself
with a slightly larger spacing to improve readability.  The ampersand
\code\& is also active.  It reads either a balanced group or
everything up to the next space, and prints the entire word (including
the ampersand) as a lambda list keyword using the
\texref{lambdalistkeyword} command.

Each definition is assigned to a category.  The name of the category is
displayed at the right-hand margin of the signature line.

You can insert multiple signatures if you want to document the
entities together.  For example,
\begin{display}
\begin{verbatim}
\begin{definition}
\defun[double]{sin}(double \var{x})
\defun[double]{cos}(double \var{x})
Return the sine or cosine of \var{x} respectively.
\end{definition}
\end{verbatim}
\end{display}
is printed as
\begin{display}
\begin{definition}
\defun[double]{sin}(double \var{x})
\defun[double]{cos}(double \var{x})
Return the sine or cosine of \var{x} respectively.
\end{definition}
\end{display}

There are three universal \symb{\def}\ldots{} commands with an
explicit category name and several specific commands with an implicit
category name.  See table~\ref{tab:def} on page~\pageref{tab:def}, for
a list of all these commands.

\begin{table}
\begin{tabular}{l l c c}
\hline
\bf Command &\bf Category &\bf Arguments &\bf Typed \\
\hline
\symb{\deffn} & \textit{explicit} & \checkmark & \checkmark \\
\symb{\defun} & \defunname & \checkmark & \checkmark \\
\symb{\defmacro} & \defmacroname & \checkmark & \checkmark \\
\symb{\defgeneric} & \defgenericname & \checkmark & \checkmark \\
\symb{\defmethod} & \defmethodname  & \checkmark & \checkmark \\
\symb{\defspec} & \defspecname  & \checkmark & \checkmark \\
\symb{\defcom} & \defcomname & \checkmark & \checkmark \\
\hline
\symb{\defvr} & \textit{explicit} & & \checkmark \\
\symb{\defvar} & \defvarname & & \checkmark \\
\symb{\defparam} & \defparamname & & \checkmark \\
\symb{\defconst} & \defconstname & & \checkmark \\
\symb{\defopt} & \defoptname & & \checkmark \\
\hline
\symb{\deftp} & \textit{explicit} & & \\
\symb{\deftype} & \deftypename & & \\
\symb{\defstruct} & \defstructname & & \\
\symb{\defclass} & \defclassname & & \\
\symb{\defcondition} & \defconditionname & & \\
\hline
\symb{\deftex} & \deftexname & \checkmark & \\
\symb{\defenv} & \defenvname & \checkmark & \\
\symb{\defcount} & \defcountname & & \\
\symb{\defdimen} & \defdimenname & & \\
\symb{\defskip} & \defskipname & & \\
\symb{\defmuskip} & \defmuskipname & & \\
\symb{\defbox} & \defboxname & & \\
\symb{\deftoks} & \deftoksname & & \\
\symb{\defif} & \defifname & & \\
\hline
\end{tabular}
\captionof{table}{Definition Commands}
\label{tab:def}
\end{table}


\subsubsection{Definition Environment}

% definition
% ----------------------------------------------------------------------
\begin{definition}
\envindex{definition}
\envanchor{definition}
\defenv{definition}
An environment for documenting commands and application programming
interfaces.  Inside a definition environment, use the \symb{\deffn},
\symb{\defvr}, or \symb{\deftp} command to insert the signature of
a function, variable, or data type like entity respectively.  The
documentation of the entity is added below the signature line.
\end{definition}

\begin{definition}
\texindex{definitionfont}
\texanchor{definitionfont}
\deftex{\definitionfont}
The main font for the signature line of a definition.

Default is ‘\code{\\normalfont\\slshape}’.  If the definition
takes arguments, mark the parameters with either \texref{var}
or \texref{meta} so that the documentation can refer them.
\end{definition}

\begin{definition}
\texindex{definitioncategoryfont}
\texanchor{definitioncategoryfont}
\deftex{\definitioncategoryfont}
The font for the category of a definition.

Default is ‘\code{\\normalfont\\bfseries\\small}’.  If the document
class is a \KOMAScript{} class, \symb{\definitioncategoryfont} uses
the respective \code{disposition} font.
\end{definition}

\begin{definition}
\texindex{lambdalistkeyword}
\texanchor{lambdalistkeyword}
\deftex{\lambdalistkeyword} \meta{word}
Format \meta{word} as a lambda list keyword.  For example,
‘\code{\\lambdalistkeyword\{\\\&rest\}}’ is printed as
\lambdalistkeyword{\&rest}.
\end{definition}

\begin{definition}
\texindex{lambdalistkeywordfont}
\texanchor{lambdalistkeywordfont}
\deftex{\lambdalistkeywordfont}
The font for a lambda list keyword.

Default is ‘\code{\\normalfont\\small}’.
\end{definition}


\subsubsection{Functions and Similar Entities}

\index{function definition}
\index{definition!function}
This section describes the commands for documenting functions and
similar entities.

% deffn
% ----------------------------------------------------------------------
\begin{definition}
\texindex{deffn}
\texanchor{deffn}
\deftex{\deffn} \meta{category} \code{*}\kleene? \meta{data-type}\kleene? \meta{symbol-name} \meta{lambda-list}
The \symb{\deffn} command is the general definition command for
functions or similar entities that may take arguments.
\begin{itemize}
\item First argument \meta{category} is the category name, for
example, “Function” if the entity is a function.
\item Optional second argument \meta{data-type} is the data type of
the return value of the function.
\item Third argument \meta{symbol-name} is the name of the function.
\item Fourth argument \meta{lambda-list} are the parameters of the
function.  The \meta{lambda-list} argument is either a balanced group
(in which case it can span multiple lines) or the rest of the line.
\end{itemize}
If the star variant is specified, \meta{data-type} is printed on a
line by itself.  Otherwise, the data type is put on the same line as
the function name and the arguments.
\end{definition}

% defun
% ----------------------------------------------------------------------
\begin{definition}
\texindex{defun}
\texanchor{defun}
\texindex{defunname}
\texanchor{defunname}
\deftex{\defun} \code{*}\kleene? \meta{data-type}\kleene? \meta{symbol-name} \meta{lambda-list}
The \symb{\defun} command is like the \texref{deffn} command with the
implicit category name “\defunname”.  The actual category name is
provided by the \symb{\defunname} command.
\end{definition}

% defmacro
% ----------------------------------------------------------------------
\begin{definition}
\texindex{defmacro}
\texanchor{defmacro}
\texindex{defmacroname}
\texanchor{defmacroname}
\deftex{\defmacro} \code{*}\kleene? \meta{data-type}\kleene? \meta{symbol-name} \meta{lambda-list}
The \symb{\defmacro} command is like the \texref{deffn} command with
the implicit category name “\defmacroname”.  The actual category name
is provided by the \symb{\defmacroname} command.
\end{definition}

% defgeneric
% ----------------------------------------------------------------------
\begin{definition}
\texindex{defgeneric}
\texanchor{defgeneric}
\texindex{defgenericname}
\texanchor{defgenericname}
\deftex{\defgeneric} \code{*}\kleene? \meta{data-type}\kleene? \meta{symbol-name} \meta{lambda-list}
The \symb{\defgeneric} command is like the \texref{deffn} command with
the implicit category name “\defgenericname”.  The actual category
name is provided by the \symb{\defgenericname} command.

Generic functions are CLOS\footnote{Common Lisp Object System}
entities describing a collection of methods with the same symbol
name and a similar argument list structure.  Other object-oriented
programming languages name this concept abstract method or interface.
To use another term, say, for example,
\begin{display}
\begin{verbatim}
\renewcommand\defgenericname{Abstract Method}
\end{verbatim}
\end{display}
at the beginning of your document.
\end{definition}

% defmethod
% ----------------------------------------------------------------------
\begin{definition}
\texindex{defmethod}
\texanchor{defmethod}
\texindex{defmethodname}
\texanchor{defmethodname}
\deftex{\defmethod} \code{*}\kleene? \meta{data-type}\kleene? \meta{symbol-name} \meta{lambda-list}
The \symb{\defmethod} command is like the \texref{deffn} command with
the implicit category name “Method”.  The actual category name is
provided by the \symb{\defmethodname} command.
\end{definition}

% defspec
% ----------------------------------------------------------------------
\begin{definition}
\texindex{defspec}
\texanchor{defspec}
\texindex{defspecname}
\texanchor{defspecname}
\deftex{\defspec} \code{*}\kleene? \meta{data-type}\kleene? \meta{symbol-name} \meta{lambda-list}
The \symb{\defspec} command is like the \texref{deffn} command with
the implicit category name “\defspecname”.  The actual category name
is provided by the \symb{\defspecname} command.

A special form is a primitive Lisp expression that looks like a
function call, but is treated specially by the Lisp system.
\end{definition}

% defcom
% ----------------------------------------------------------------------
\begin{definition}
\texindex{defcom}
\texanchor{defcom}
\texindex{defcomname}
\texanchor{defcomname}
\deftex{\defcom} \code{*}\kleene? \meta{data-type}\kleene? \meta{symbol-name} \meta{lambda-list}
The \symb{\defcom} command is like the \texref{deffn} command with the
implicit category name “\defcomname”.  The actual category name is
provided by the \symb{\defcomname} command.
\end{definition}


\subsubsection{Variables and Similar Entities}

\index{variable definition}
\index{definition!variable}
This section describes the commands for documenting variables and
similar entities.

% defvr
% ----------------------------------------------------------------------
\begin{definition}
\texindex{defvr}
\texanchor{defvr}
\deftex{\defvr} \meta{category} \code{*}\kleene? \meta{data-type}\kleene? \meta{symbol-name}
The \symb{\defvr} command is the general definition command for
variables or similar entities.
\begin{itemize}
\item First argument \meta{category} is the category name, for
example, “Variable” if the entity is a variable.
\item Optional second argument \meta{data-type} is the data type of
the variable.
\item Third argument \meta{symbol-name} is the name of the variable.
\end{itemize}
If the star variant is specified, \meta{data-type} is printed on a
line by itself.  Otherwise, the data type is put on the same line as
the variable name.
\end{definition}

For example,
\begin{display}
\begin{verbatim}
\begin{definition}
\defvr{Variable}*[char *]{program_invocation_name}
This variable’s value is the name that was used to invoke
the program running in the current process.
\end{definition}
\end{verbatim}
\end{display}
is printed as
\begin{display}
\begin{definition}
\defvr{Variable}*[char *]{program_invocation_name}
This variable’s value is the name that was used to invoke
the program running in the current process.
\end{definition}
\end{display}

There are specialized commands for common category names.

% defvar
% ----------------------------------------------------------------------
\begin{definition}
\texindex{defvar}
\texanchor{defvar}
\texindex{defvarname}
\texanchor{defvarname}
\deftex{\defvar} \code{*}\kleene? \meta{data-type}\kleene? \meta{symbol-name}
The \symb{\defvar} command is like the \texref{defvr} command with the
implicit category name “\defvarname”.  The actual category name is
provided by the \symb{\defvarname} command.
\end{definition}

% defparam
% ----------------------------------------------------------------------
\begin{definition}
\texindex{defparam}
\texanchor{defparam}
\texindex{defparamname}
\texanchor{defparamname}
\deftex{\defparam} \code{*}\kleene? \meta{data-type}\kleene? \meta{symbol-name}
The \symb{\defparam} command is like the \texref{defvr} command with
the implicit category name “\defparamname”.  The actual category name
is provided by the \symb{\defparamname} command.
\end{definition}

% defconst
% ----------------------------------------------------------------------
\begin{definition}
\texindex{defconst}
\texanchor{defconst}
\texindex{defconstname}
\texanchor{defconstname}
\deftex{\defconst} \code{*}\kleene? \meta{data-type}\kleene? \meta{symbol-name}
The \symb{\defconst} command is like the \texref{defvr} command with
the implicit category name “\defconstname”.  The actual category name
is provided by the \symb{\defconstname} command.
\end{definition}

% defopt
% ----------------------------------------------------------------------
\begin{definition}
\texindex{defopt}
\texanchor{defopt}
\texindex{defoptname}
\texanchor{defoptname}
\deftex{\defopt} \code{*}\kleene? \meta{data-type}\kleene? \meta{symbol-name}
The \symb{\defopt} command is like the \texref{defvr} command with the
implicit category name “\defoptname”.  The actual category name is
provided by the \symb{\defoptname} command.
\end{definition}


\subsubsection{Data Types and Similar Entities}

\index{data type definition}
\index{definition!data type}
This section describes the commands for documenting data types and
similar entities.

% deftp
% ----------------------------------------------------------------------
\begin{definition}
\texindex{deftp}
\texanchor{deftp}
\deftex{\deftp} \meta{category} \meta{symbol-name}
The \symb{\deftp} command is the general definition command for data
types or similar entities.
\begin{itemize}
\item First argument \meta{category} is the category name, for
example, “Data Type” if the entity is a data type.
\item Second argument \meta{symbol-name} is the name of the data type.
\end{itemize}
\end{definition}

% deftype
% ----------------------------------------------------------------------
\begin{definition}
\texindex{deftype}
\texanchor{deftype}
\texindex{deftypename}
\texanchor{deftypename}
\deftex{\deftype} \meta{symbol-name}
The \symb{\deftype} command is like the \texref{deftp} command with
the implicit category name “\deftypename”.  The actual category name
is provided by the \symb{\deftypename} command.
\end{definition}

% defstruct
% ----------------------------------------------------------------------
\begin{definition}
\texindex{defstruct}
\texanchor{defstruct}
\texindex{defstructname}
\texanchor{defstructname}
\deftex{\defstruct} \meta{symbol-name}
The \symb{\defstruct} command is like the \texref{deftp} command with
the implicit category name “\defstructname”.  The actual category name
is provided by the \symb{\defstructname} command.
\end{definition}

% defclass
% ----------------------------------------------------------------------
\begin{definition}
\texindex{defclass}
\texanchor{defclass}
\texindex{defclassname}
\texanchor{defclassname}
\deftex{\defclass} \meta{symbol-name}
The \symb{\defclass} command is like the \texref{deftp} command with
the implicit category name “\defclassname”.  The actual category name
is provided by the \symb{\defclassname} command.
\end{definition}

% defcondition
% ----------------------------------------------------------------------
\begin{definition}
\texindex{defcondition}
\texanchor{defcondition}
\texindex{defconditionname}
\texanchor{defconditionname}
\deftex{\defcondition} \meta{symbol-name}
The \symb{\defcondition} command is like the \texref{deftp} command
with the implicit category name “\defconditionname”.  The actual
category name is provided by the \symb{\defconditionname} command.
\end{definition}

\index{exception}
If your programming language only provides exceptions, say, for
example,
\begin{display}
\begin{verbatim}
\renewcommand\defconditionname{Exception}
\end{verbatim}
\end{display}
at the beginning of your document.


\subsubsection{\TeX{} and \LaTeX{} Entities}

There are specialized commands for documenting \TeX{} and \LaTeX{}
entities.

% deftex
% ----------------------------------------------------------------------
\begin{definition}
\texindex{deftex}
\texanchor{deftex}
\texindex{deftexname}
\texanchor{deftexname}
\deftex{\deftex} \meta{symbol-name} \meta{lambda-list}
The \symb{\deftex} command is like the \texref{deffn} command with the
implicit category name “\deftexname”.  The actual category name is
provided by the \symb{\deftexname} command.
\end{definition}

% defenv
% ----------------------------------------------------------------------
\begin{definition}
\texindex{defenv}
\texanchor{defenv}
\texindex{defenvname}
\texanchor{defenvname}
\deftex{\defenv} \meta{symbol-name} \meta{lambda-list}
The \symb{\defenv} command is like the \texref{deffn} command with the
implicit category name “\defenvname”.  The actual category name is
provided by the \symb{\defenvname} command.
\end{definition}

% defcount
% ----------------------------------------------------------------------
\begin{definition}
\texindex{defcount}
\texanchor{defcount}
\texindex{defcountname}
\texanchor{defcountname}
\deftex{\defcount} \meta{symbol-name}
The \symb{\defcount} command is like the \texref{defvr} command with
the implicit category name “\defcountname”.  The actual category name
is provided by the \symb{\defcountname} command.
\end{definition}

% defdimen
% ----------------------------------------------------------------------
\begin{definition}
\texindex{defdimen}
\texanchor{defdimen}
\texindex{defdimenname}
\texanchor{defdimenname}
\deftex{\defdimen} \meta{symbol-name}
The \symb{\defdimen} command is like the \texref{defvr} command with
the implicit category name “\defdimenname”.  The actual category name
is provided by the \symb{\defdimenname} command.
\end{definition}

% defskip
% ----------------------------------------------------------------------
\begin{definition}
\texindex{defskip}
\texanchor{defskip}
\texindex{defskipname}
\texanchor{defskipname}
\deftex{\defskip} \meta{symbol-name}
The \symb{\defskip} command is like the \texref{defvr} command with
the implicit category name “\defskipname”.  The actual category name
is provided by the \symb{\defskipname} command.
\end{definition}

% defmuskip
% ----------------------------------------------------------------------
\begin{definition}
\texindex{defmuskip}
\texanchor{defmuskip}
\texindex{defmuskipname}
\texanchor{defmuskipname}
\deftex{\defmuskip} \meta{symbol-name}
The \symb{\defmuskip} command is like the \texref{defvr} command with
the implicit category name “\defmuskipname”.  The actual category name
is provided by the \symb{\defmuskipname} command.
\end{definition}

% defbox
% ----------------------------------------------------------------------
\begin{definition}
\texindex{defbox}
\texanchor{defbox}
\texindex{defboxname}
\texanchor{defboxname}
\deftex{\defbox} \meta{symbol-name}
The \symb{\defbox} command is like the \texref{defvr} command with the
implicit category name “\defboxname”.  The actual category name is
provided by the \symb{\defboxname} command.
\end{definition}

% deftoks
% ----------------------------------------------------------------------
\begin{definition}
\texindex{deftoks}
\texanchor{deftoks}
\texindex{deftoksname}
\texanchor{deftoksname}
\deftex{\deftoks} \meta{symbol-name}
The \symb{\deftoks} command is like the \texref{defvr} command with
the implicit category name “\deftoksname”.  The actual category name
is provided by the \symb{\deftoksname} command.
\end{definition}

% defif
% ----------------------------------------------------------------------
\begin{definition}
\texindex{defif}
\texanchor{defif}
\texindex{defifname}
\texanchor{defifname}
\deftex{\defif} \meta{symbol-name}
The \symb{\defif} command is like the \texref{defvr} command with the
implicit category name “\defifname”.  The actual category name is
provided by the \symb{\defifname} command.
\end{definition}


\subsection{Date and Time}

\index{date}
\index{time}
\index{ISO 8601}
\styindex{rsdatetime}
The \code{rsdatetime} package provides macros to print the date and
time in ISO~8601 format.  This is a lightweight solution.  If you need
more advanced capabilities, use the
\href{https://mirrors.ctan.org/macros/latex/contrib/datetime2/datetime2.pdf}{\code{datetime2}}
package.

% isodate
% ----------------------------------------------------------------------
\begin{definition}
\texindex{isodate}
\texanchor{isodate}
\deftex{\isodate} \code{*}\kleene? \meta{precision}\kleene?
Insert the current date in extended ISO 8601 format.
\begin{itemize}
\item Optional argument \meta{precision} specifies the precision.
Value is either \code{year}, \code{month}, \code{day}, \code{hour},
\code{minute}, or \code{second}.  Default is \code{day}.  If the
precision is a time of day unit, the time is separated from the date
by a ‘T’ character.
\end{itemize}
If the star indicator is present, utilize the basic format.
\end{definition}

% isotime
% ----------------------------------------------------------------------
\begin{definition}
\texindex{isotime}
\texanchor{isotime}
\deftex{\isotime} \code{*}\kleene? \meta{precision}\kleene?
Insert the current time in extended ISO 8601 format.
\begin{itemize}
\item Optional argument \meta{precision} specifies the precision.
Value is either \code{hour}, \code{minute}, or \code{second}.  Default
is \code{minute}.
\end{itemize}
If the star indicator is present, utilize the basic format.
\end{definition}

Whether or not a second precision is actually supported depends on the
underlying \TeX{} engine.  If not, the second of the minute defaults
to zero.

Table~\ref{tab:date} on page~\pageref{tab:date} lists the date and
time formats produced by the \texref{isodate} and \texref{isotime}
commands.

\begingroup
% Example from the ISO 8601 standard: 1985-04-12T10:15:30
\year=1985\month=4\day=12\time=615
\begin{display}*
\begin{tabular}{l l l l}
\hline
\multicolumn{2}{c}{Extended Format} &
\multicolumn{2}{c}{Basic Format} \\
\hline
\hfil Command  & \hfil Result &
\hfil Command  & \hfil Result \\
\hline
\code{\\isodate}          & \isodate          &
\code{\\isodate*}         & \isodate*         \\
\code{\\isodate[year]}    & \isodate[year]    &
\code{\\isodate*[year]}   & \isodate*[year]   \\
\code{\\isodate[month]}   & \isodate[month]   &
\code{\\isodate*[month]}  & \isodate*[month]  \\
\code{\\isodate[day]}     & \isodate[day]     &
\code{\\isodate*[day]}    & \isodate*[day]    \\
\code{\\isodate[hour]}    & \isodate[hour]    &
\code{\\isodate*[hour]}   & \isodate*[hour]   \\
\code{\\isodate[minute]}  & \isodate[minute]  &
\code{\\isodate*[minute]} & \isodate*[minute] \\
\hline
\code{\\isotime}          & \isotime          &
\code{\\isotime*}         & \isotime*         \\
\code{\\isotime[hour]}    & \isotime[hour]    &
\code{\\isotime*[hour]}   & \isotime*[hour]   \\
\code{\\isotime[minute]}  & \isotime[minute]  &
\code{\\isotime*[minute]} & \isotime*[minute] \\
\hline
\end{tabular}
\captionof{table}{Date and Time Formats}
\label{tab:date}
\end{display}
\endgroup


\clearpage
\section{Index}

\subsection*{Concept Index}
\printindex

\subsection*{Index of \LaTeX{} Macros}
\printindex[sym]
\end{document}

%% rsuserguide.tex ends here
